% =============================================================================
% Chapter 5: Experimental Setup
% Source of truth: Thesis_Final_Handover.md (6 October 2026) and
% spi_results_data.json. Every number is taken from those two files unchanged.
% Tables and figure blocks are copied verbatim from tables.tex and
% figure_captions.tex so that they cannot drift from the generated versions.
% =============================================================================

\chapter{Experimental Setup}
\label{ch:experiment}

The Causal Chain Derivation Method of Chapter~\ref{ch:methodology} yields a
register of nine indicators. Whether that register is usable cannot be settled by
argument alone, so this chapter describes the simulation campaigns against which
it is evaluated and the measures by which the evaluation is judged. One
unprotected left turn exposes three conflict geometries under controlled and
reproducible conditions, which is what allows the scalability question of
Section~\ref{sec:intro_objective} to be answered within a single manoeuvre.

Section~\ref{sec:exp_scenario} describes the manoeuvre and its three conflicts.
Section~\ref{sec:exp_environment} states the ego configuration and the simulation
parameters. Section~\ref{sec:exp_parameters} gives the quantities sampled per
run. Section~\ref{sec:exp_audit} accounts for the difference between the rows
recorded and the runs analysed. Section~\ref{sec:exp_collision} fixes what counts
as a collision, which for the pedestrian chain is not self-evident.
Section~\ref{sec:exp_criteria} defines the evaluation measures, and
Section~\ref{sec:exp_calibration} gives the calibrated metric thresholds.

% =============================================================================
\section{Scenario and Conflict Geometries}
\label{sec:exp_scenario}
\label{sec:exp_manoeuvre}

A single manoeuvre is used throughout, because three conflicts of different
character arise within it and can therefore be compared without changing the
driving task. The ego vehicle enters the intersection from the south and turns
left across the opposing lane. All traffic signals are frozen green, so no signal
resolves any conflict on the ego vehicle's behalf and every conflict has to be
resolved by the vehicle itself. Figure~\ref{fig:5_1_scenario} shows the
intersection and the three conflict geometries.

% ---------------------------------------------------------------------------
% Figure 5.1: schematic conflict geometry (Graphics/leftturn_conflict_geometry.pdf).
\begin{figure}[htbp]
    \centering
    \includegraphics[width=16cm]{leftturn_conflict_geometry.pdf}
    \caption{Schematic of the unprotected left turn on Town10HD. The ego vehicle enters from
    the south and turns left under a frozen green signal; oncoming vehicles cross
    the turn arc at an angle, pedestrians cross the leg the ego vehicle turns
    into, and a following vehicle closes along the ego vehicle's path.}
    \label{fig:5_1_scenario}
\end{figure}

The three conflicts differ in the way they close. In the oncoming-vehicle chain
(ONC) the partner crosses the turn arc at an angle, which is the geometry
classified as left turn across path, opposite direction. In the pedestrian chain
(PED) walkers cross the leg that the ego vehicle turns into, late in the turn, so
the conflict closes laterally. In the following-vehicle chain (REAR) a follower
closes along the ego vehicle's own path whenever the ego vehicle brakes, so the
conflict closes longitudinally. One manoeuvre therefore carries a crossing, a
lateral and a following conflict at once.

% =============================================================================
\section{Ego Vehicle and Simulation Environment}
\label{sec:exp_environment}

The quantity under evaluation is the indicator, not a driving function, and the
ego configuration is chosen accordingly: it has to produce conflicts of all three
kinds reliably and reproducibly, not to resolve them well.

The ego vehicle is driven by the CARLA Traffic Manager with an added
pedestrian-stop rule and without gap acceptance. It therefore enters the turn
without evaluating whether the gap in oncoming traffic is sufficient. The
evaluation consequently tests the indicators on trajectories rather than on a
driving function, and the oncoming conflicts that arise are the result of a
missing capability rather than of a functional insufficiency in the sense of
ISO~21448 \cite{iso21448}.

Simulation runs in CARLA~0.9.15 \cite{dosovitskiy2017carla} on the map Town10HD,
in synchronous mode at 20\,Hz, that is at a fixed step of 0.05\,s. Each run
covers 400 steps, or 20\,s, and executes in an isolated subprocess under one
seed, so that a failure in one run cannot affect another. Only raw telemetry is
stored: the ego vehicle's state and control inputs, the pose and velocity of
every actor, and every collision event with the identity of the actor involved.
No indicator is computed during the run. All indicators are computed offline from
the stored ground truth, which keeps the derivation independent of the simulated
vehicle's own perception and makes every value reproducible from the frozen data.
The datasets were frozen on 28~September 2026.

Three campaigns were executed. They differ in the scenario generator as well as
in the controller configuration, so a comparison between them is a comparison of
two known groups rather than a controlled variation of one factor.
Table~\ref{tab:campaigns} gives their composition.

% ---------------------------------------------------------------------------
\begin{table}[htbp]
    \centering
    \caption{Simulation campaigns. The table gives the configuration and the
    number of analysed runs of each campaign.}
    \label{tab:campaigns}
    \small
    \begin{tabular}{@{}llr@{}}
        \toprule
        {\bfseries Campaign} & {\bfseries Configuration} & {\bfseries Analysed runs} \\
        \midrule
        Campaign 1 & collisions provoked, high-collision configuration & 19,307 \\
        Campaign 2 & realistic driving, low-collision configuration & 33,797 \\
        Campaign 3 & realistic driving, low-collision configuration & 41,712 \\
        \midrule
        Total      &                                                 & 94,816 \\
        \bottomrule
    \end{tabular}
    \par\smallskip\footnotesize\raggedright Campaign~3 is statistically
    indistinguishable from Campaign~2; see Section~\ref{sec:results_poolability}.
\end{table}

% =============================================================================
\section{Sampled Parameters}
\label{sec:exp_parameters}
\label{sec:exp_pipeline}

Reproducibility requires that the variation between runs be explicit. Every run
draws its actors and its environment from the ranges given in
Tables~\ref{tab:actor_parameters} and~\ref{tab:environment_parameters} under the
run's seed, so that any run can be regenerated exactly.

% ---------------------------------------------------------------------------
\begin{table}[htbp]
    \centering
    \caption{Actor parameters sampled per run. Each run draws the number,
    appearance and behaviour of its conflict partners from the ranges given.}
    \label{tab:actor_parameters}
    \small
    \begin{tabular}{@{}lll@{}}
        \toprule
        {\bfseries Actor} & {\bfseries Parameter} & {\bfseries Range or set} \\
        \midrule
        Oncoming vehicle & number            & 1 to 3 \\
                         & vehicle model     & four models \\
                         & speed             & 7 to 12\,m/s \\
                         & turn direction    & straight, left, right \\
                         & obeys signal      & yes, no \\
                         & yields            & yes, no \\
                         & optional braking  & 4.0\,m/s\textsuperscript{2} \\
                         & launch            & staggered \\
        \midrule
        Pedestrian       & number            & 1 to 3 \\
                         & walker model      & four models \\
                         & crossing route    & eight routes over four crosswalks \\
                         & speed             & 0.9 to 1.7\,m/s \\
                         & group crossing    & probability 0.5 \\
        \midrule
        Following vehicle & vehicle model    & four models \\
                         & initial gap       & 4 to 20\,m \\
                         & speed             & 8 to 12\,m/s \\
                         & reaction delay    & 0 to 1.5\,s \\
                         & minimum gap       & 2 to 6\,m \\
        \bottomrule
    \end{tabular}
\end{table}

% ---------------------------------------------------------------------------
\begin{table}[htbp]
    \centering
    \caption{Environment parameters sampled per run. Cloud cover, precipitation
    and deposits are drawn together, so wet road surfaces accompany rain.}
    \label{tab:environment_parameters}
    \small
    \begin{tabular}{@{}ll@{}}
        \toprule
        {\bfseries Parameter} & {\bfseries Range} \\
        \midrule
        Cloud cover, precipitation, deposits & 0 to 100\,\%, coupled \\
        Wind                                 & sampled \\
        Fog density                          & sampled \\
        Wetness                              & sampled \\
        Sun altitude                         & $-20$ to 90\textdegree \\
        Fog visibility                       & 10 to 200\,m \\
        \bottomrule
    \end{tabular}
\end{table}

% =============================================================================
\section{Data Audit}
\label{sec:exp_audit}

The number of rows written by the simulation campaigns is not the number of runs
available for analysis, and the difference is large enough that it has to be
accounted for rather than assumed away. A failed attempt was written as a row and
then retried under the same seed, so rows overstate the number of runs.
Figure~\ref{fig:5_1_data_audit} traces the reduction from recorded rows to
analysed runs.

% ---------------------------------------------------------------------------
% Copied verbatim from figure_captions.tex
\begin{figure}[htbp]
\centering
\includegraphics[width=16cm]{fig_5_1_data_audit.pdf}
\caption{Data audit of the frozen simulation campaigns. Failed attempts were logged as rows and retried under the same seed; keeping one complete run per seed and excluding the first 500 rows of Campaign 1 leaves 94,816 analysed runs.}
\label{fig:5_1_data_audit}
\end{figure}

Of 113,787 registered rows, 95,586 carry a distinct seed. Of those, 95,143 are
complete runs that reached all 400 of 400 steps. The first 500 rows of
Campaign~1, corresponding to 327 complete runs, are excluded because they were
produced while the CARLA Traffic Manager parameters were still being adjusted to
the requirements of the scenario, so they do not belong to the final
configuration. This leaves 94,816 analysed runs, which is the population behind
every result in Chapter~\ref{ch:results}.

% =============================================================================
\section{Collision Definition}
\label{sec:exp_collision}

The reference event of the whole evaluation is the collision, so its definition
decides what every leading indicator is scored against. For the two vehicle
chains the definition is immediate: any recorded contact with the oncoming
vehicle or with the following vehicle counts. For the pedestrian chain it is not,
because CARLA walkers follow a planned path and do not avoid vehicles. A walker
whose route crosses a stationary ego vehicle walks into it, and the simulator
emits a contact event that is indistinguishable in the log from a contact the ego
vehicle caused.

A pedestrian contact therefore counts as a collision only if the ego vehicle
moves at at least 0.3\,m/s at the simulation step before the contact event.
Contacts with a standing ego vehicle are recorded but not counted.
Figure~\ref{fig:5_2_pedestrian_contacts} and Table~\ref{tab:pedestrian_contacts}
give the audit on which that rule rests.

% ---------------------------------------------------------------------------
% Copied verbatim from figure_captions.tex
\begin{figure}[htbp]
\centering
\includegraphics[width=16cm]{fig_5_2_pedestrian_contacts.pdf}
\caption{Pedestrian contacts by type and campaign. Each bar classifies the first pedestrian contact of every contact run by the motion one simulation step before the contact; only contacts while the ego moves at least 0.3\,m/s count as pedestrian collisions.}
\label{fig:5_2_pedestrian_contacts}
\end{figure}

% ---------------------------------------------------------------------------
\begin{table}[htbp]
    \centering
    \caption{Audit of the pedestrian contacts. The table classifies the first
    pedestrian contact of each of the 1,323 contact runs by the motion of the ego
    vehicle one simulation step before the contact.}
    \label{tab:pedestrian_contacts}
    \small
    \begin{tabular}{@{}lrrrr@{}}
        \toprule
        {\bfseries Campaign} &
        {\bfseries \begin{tabular}[b]{@{}c@{}}Ego drives into\\the pedestrian\end{tabular}} &
        {\bfseries \begin{tabular}[b]{@{}c@{}}Pedestrian walks into\\the moving ego\end{tabular}} &
        {\bfseries \begin{tabular}[b]{@{}c@{}}Pedestrian walks into\\the standing ego\end{tabular}} &
        {\bfseries Contacts} \\
        \midrule
        Campaign 1 & 812 & 2 & 28  & 842 \\
        Campaign 2 & 0   & 1 & 193 & 194 \\
        Campaign 3 & 0   & 0 & 287 & 287 \\
        \midrule
        Total      & 812 & 3 & 508 & 1,323 \\
        \bottomrule
    \end{tabular}
    \par\smallskip\footnotesize\raggedright The first two columns count as
    pedestrian collisions, the third does not.
\end{table}

The two populations differ in kind and not only in label. Contacts in which the
ego vehicle drives into the pedestrian occur at a median ego speed of 7.1\,m/s
one step before contact, and the pedestrian is in front of the ego vehicle in 800
of 812 cases. Contacts with a standing ego vehicle occur at an ego speed of
0.0\,m/s and carry a median contact impulse of 34\,N\,s, against 310\,N\,s for
contacts in which the ego vehicle is driving, which is about one tenth.

Applying the rule leaves 815 pedestrian collisions: 814 in Campaign~1, one in
Campaign~2 and none in Campaign~3. The effect is concentrated in the two
low-collision campaigns, where 480 of the 481 recorded contacts occur with a
standing ego vehicle.

% =============================================================================
\section{Evaluation Measures}
\label{sec:exp_criteria}

An indicator is evaluated by how its violations relate to the collisions of its
own chain, so each eligible run is classified before any measure is computed. A
run of a given chain is a true positive (TP) if the indicator is violated before
a collision of that chain, a false positive (FP) if it is violated and no
collision of that chain occurs, a false negative (FN) if a collision occurs
without an earlier violation, and a true negative (TN) if neither happens. Only
violations that precede the collision are counted, so that a violation can never
be a consequence of the event it is meant to anticipate.

Three measures follow directly from that classification. Precision, in
Equation~(\ref{eq:exp_precision}), is the share of violating runs that end in a
collision of the chain.

\begin{equation}
    \mathrm{Precision} = \frac{\mathrm{TP}}{\mathrm{TP} + \mathrm{FP}}
    \label{eq:exp_precision}
\end{equation}
\equationentry{Precision}

Recall, in Equation~(\ref{eq:exp_recall}), is the share of collision runs
preceded by a violation.

\begin{equation}
    \mathrm{Recall} = \frac{\mathrm{TP}}{\mathrm{TP} + \mathrm{FN}}
    \label{eq:exp_recall}
\end{equation}
\equationentry{Recall}

The false-alarm rate, in Equation~(\ref{eq:exp_falsealarm}), is the share of
collision-free eligible runs that nevertheless record a violation.

\begin{equation}
    \mathrm{False\text{-}alarm\ rate} = \frac{\mathrm{FP}}{\mathrm{FP} + \mathrm{TN}}
    \label{eq:exp_falsealarm}
\end{equation}
\equationentry{False-alarm rate}

Precision and recall describe the indicator but do not state how much a violation
changes the collision probability. The relative risk, in
Equation~(\ref{eq:exp_relrisk}), gives that directly as the ratio of the
collision probability among violating runs to the collision probability among
runs without a violation.

\begin{equation}
    \mathrm{RR} = \frac{\mathrm{TP} / (\mathrm{TP} + \mathrm{FP})}
                       {\mathrm{FN} / (\mathrm{FN} + \mathrm{TN})}
    \label{eq:exp_relrisk}
\end{equation}
\equationentry{Relative risk of a collision after a violation}

\textbf{Variable key:} TP, number of eligible runs with a violation before a
collision of the chain, dimensionless; FP, number of eligible runs with a
violation and without a collision of the chain, dimensionless; FN, number of
eligible runs with a collision of the chain and without an earlier violation,
dimensionless; TN, number of eligible runs with neither, dimensionless; RR,
relative risk, dimensionless.

Every proportion is reported with a Wilson score interval and every relative risk
with the Katz log interval of Equation~(\ref{eq:exp_katz}), both at 95\,\%
[CITATION NEEDED: Wilson score interval] [CITATION NEEDED: Katz log interval for
the relative risk].

\begin{equation}
    \ln \mathrm{RR} \pm z \sqrt{\frac{1}{\mathrm{TP}}
        - \frac{1}{\mathrm{TP} + \mathrm{FP}}
        + \frac{1}{\mathrm{FN}}
        - \frac{1}{\mathrm{FN} + \mathrm{TN}}}
    \label{eq:exp_katz}
\end{equation}
\equationentry{Katz log interval of the relative risk}

\textbf{Variable key:} $z$, standard normal quantile, here 1.96 for a two-sided
95\,\% interval, dimensionless.

Where a cell of the classification is zero the interval is not defined, so 0.5 is
added to every cell before the interval is computed [CITATION NEEDED:
Haldane-Anscombe correction]. Where FN is zero the corrected relative risk is
reported as a lower bound rather than as a point estimate.

The SPI threshold T is defined on the violation count of a monitoring window
rather than on a single run. For a window of eligible runs, T is the smallest
count at which the observed number of violations becomes improbable under the
violation rate of the realistic-driving campaigns, as given in
Equation~(\ref{eq:exp_spithreshold}).

\begin{equation}
    T = \min \left\{ k \in \mathbb{N} \;:\;
        P(X \geq k) \leq \alpha,\quad X \sim \mathrm{Bin}(n, p_0) \right\}
    \label{eq:exp_spithreshold}
\end{equation}
\equationentry{SPI threshold as an exact one-sided binomial limit}

\textbf{Variable key:} $k$, violation count of the window, dimensionless; $n$,
number of eligible runs in the window, dimensionless; $p_0$, violation rate of
Campaigns~2 and~3 for the indicator, dimensionless; $\alpha$, significance level,
here 0.01, dimensionless.

Monitoring windows hold 200 consecutive runs. The window is the unit over which
the rate is checked; it is not the unit of T, which is stated in the general SPI
form of Section~\ref{sec:method_thresholds}. Because T is an exact limit at
$\alpha = 0.01$, about one window in a hundred exceeds it in the campaigns from
which it is derived, by construction.

% =============================================================================
\section{Calibration of the Metric Thresholds}
\label{sec:exp_calibration}

The metric threshold $\theta$ is calibrated once and then applied unchanged, so
that no indicator is tuned on the collisions against which it is later scored.
The calibration population is the set of uneventful runs of Campaign~1, that is
runs that end without a collision and without an emergency braking event. One
value per run enters the calibration, namely the minimum or the maximum of the
metric over the run, because consecutive simulation steps are strongly
autocorrelated and would otherwise be counted as independent observations.
Table~\ref{tab:metric_thresholds} gives the result.

% ---------------------------------------------------------------------------
% Copied verbatim from tables.tex
\begin{table}[htbp]
\centering
\caption{Metric thresholds. Each metric threshold is a percentile of the Campaign~1 runs without collision and without emergency braking.}
\label{tab:metric_thresholds}
\small
\begin{tabular}{@{}llllll@{}}
\toprule
{\bfseries Indicator} & {\bfseries Chain} & {\bfseries Violation} & {\bfseries $\theta$} & {\bfseries 95\,\% interval} & {\bfseries Calibration runs} \\
\midrule
GAP\_TTC & ONC & below & 0.913\,s & 0.905 to 0.920\,s & 13,290 \\
REQ\_DECEL & ONC & above & 7.549\,m/s\textsuperscript{2} & 7.253 to 7.903\,m/s\textsuperscript{2} & 13,290 \\
PED\_CPA & PED & below & 1.174\,s & 1.138 to 1.261\,s & 2,070 \\
REAR\_TTC & REAR & below & 2.679\,s & 2.596 to 2.794\,s & 12,505 \\
\bottomrule
\end{tabular}
\par\smallskip\footnotesize\raggedright 10th percentile for time-based metrics, 95th percentile for REQ\_DECEL; bootstrap percentile intervals.
\end{table}

The bootstrap intervals are narrow relative to the thresholds themselves, so the
percentile is well determined by the calibration sample. The interval of
PED\_CPA is the widest, at 1.138 to 1.261\,s, and rests on 2,070 calibration
runs, against 13,290 for the two oncoming-vehicle indicators.

The thresholds are not transferable between configurations without checking. The
metric threshold of REQ\_DECEL shifts by 21.5\,\% and that of PED\_CPA by
37.9\,\% between Campaign~1 and the realistic-driving campaigns, which is why
calibration is performed once on Campaign~1 and the resulting value is applied
unchanged rather than refitted per campaign. Section~\ref{sec:results_poolability}
reports the shift for every indicator.

Two reaction-tier indicators carry no metric threshold. AEB\_ONC and AEB\_PED
fire on an event, namely a brake command of at least 0.8 together with a
line-of-sight time to collision of at most 1.5\,s to the gated partner of the
chain in the same simulation step. Their calibration population of 16,950 runs is
used only to establish that the event is rare under uneventful driving.